\chapter{Sviluppo}

\section{Testing automatizzato}

Il testing automatizzato e' stato realizzato con JUnit ed eseguito tramite il
task Gradle \texttt{test}, configurato con \texttt{useJUnitPlatform()}. I test
sono concentrati sulla parte di modello dell'applicazione. In particolare
sono stati testati i componenti che gestiscono giocatore, mano,
mazzo, armata, creature e caricamento delle definizioni delle carte. Sono
inoltre presenti test per il \texttt{BoardGame} e per le componenti
dell'intelligenza artificiale, come generazione delle azioni legali,
transizioni di stato simulato, valutazione euristica e scelta delle mosse.
In questo modo le principali regole della partita e il comportamento automatico
dell'avversario possono essere controllati in modo ripetibile a ogni esecuzione
della suite.

\section{Note di sviluppo}

\subsection{Massimiliano Ria}

\subsection{Simone Marsili}

\subsection{Francesco Dama}

\subsection{Lorenzo Mercuriali}
